\thispagestyle{empty}
\begin{latin}
\sloppy
\hbadness=2500
{\bfseries\fontsize{20pt}{28pt}\selectfont Abstract\par}
\vspace{1.1cm}

Operating a microgrid means deciding, at each time step, how to serve the load, whether to charge or discharge the battery, and how much power to exchange with the utility grid. Those decisions have to respect power balance and storage limits. A neural network trained only on examples may fit the labels and still violate the same limits. Classical optimization handles the constraints more directly, but repeatedly solving the program is not always cheap.

This project uses the idea of physics-informed neural networks: residuals of power balance and battery state-of-charge are added to the training loss. The list of operating variables follows grid-connected energy-management models. Game-theoretic profit sharing among several sites is reviewed but is not part of the numerical test.

The case study is a small grid-connected microgrid with photovoltaic generation, a battery and a local load. A feed-forward network was programmed in Python. The data come from a twenty-one-day hourly simulator with a fixed seed. On the test hours the physics-informed model reduced the power-balance residual and the bound-violation rate compared with a data-only network, while the mean squared error increased slightly. The observation is limited to this synthetic set.

\vspace{0.6cm}
\noindent\textbf{Keywords:} \KeywordsEn
\end{latin}
